\subsection{Why Traditional NLP Still Matters}
\label{sec:why-traditional-nlp}
\slidesrc{Terminology introduction: 9--10}

If LLMs can do almost everything, why study tokenisers, taggers and grammars? Because
traditional NLP methods and approaches are still required in three situations.

\begin{nbfact}{When traditional NLP is required}
\begin{enumerate}[nosep]
  \item \textbf{Not enough data}: when the amount of training data is insufficient to
        apply machine learning or deep learning methods successfully, or to train LLMs
        (e.g.\ low-resource languages, very specialised domains).
  \item \textbf{Pre- and post-processing}: when NLP pipelines need preprocessing (e.g.\
        tokenisation) or postprocessing and transformation of their output (e.g.\
        extracting knowledge from syntactic parses). This includes \emph{refining and
        validating the results} of ML, DL and LLM models.
  \item \textbf{LLMs cannot be used}, for technical reasons (cost, latency, hardware),
        legal reasons (privacy, data that cannot leave the organisation) or ethical
        reasons (lack of transparency, bias).
\end{enumerate}
\end{nbfact}

\begin{nbintuition}
Even an LLM starts with a tokeniser, a very traditional NLP component. And when an LLM
produces an answer that a program must use (a database query, a logical fact), checking
that the answer is well formed and consistent is again a job for traditional methods.
The case study below shows exactly this.
\end{nbintuition}

The course will therefore discuss (some of) the common problems of NLP following a
\emph{standard linguistic analysis approach} (Section~\ref{sec:linguistic-analysis}),
with examples of (some of) the common techniques and their common implementations.
